\chapter{Small Clause Forms: Naming, Witnesses, and Simplification}
\label{sec:p12-opt-small-clauses}

This chapter develops Sections 4--8 of Nonnengart and Weidenbach's
\emph{Computing Small Clause Normal Forms}, in its retained 2001
presentation \cite[pp.~347--365]{NonnengartWeidenbach2001SmallClauses}.
The cost being optimized is initially the number of clauses produced
before redundancy removal. It is neither the number of models nor a
prediction of a prover's running time. Formula naming, witness selection,
and simplification affect these quantities differently.

\section{Semantic interface}
\begin{hbdefinition}{def:p12-opt-interface}{Inputs and fresh symbols}
Use classical first-order logic over a declared signature \(\Sigma\),
with nonempty domains and total functions. Formulas are finite syntax
trees; an occurrence is a position in that tree. Substitution is
capture-avoiding. Rectification makes distinct binders use distinct
variables, also disjoint from free variables.
For arbitrary open input \(F\), existential-assignment satisfiability is
handled by first taking \(\exists\operatorname{FV}(F)\,F\).
If instead the problem requires universal closure, that closure is
part of the input specification. These two problems are different.

All newly introduced predicates and functions are outside \(\Sigma\).
An expansion preserves the domain and every interpretation in \(\Sigma\);
a reduct forgets only the added symbols. Choosing a name merely absent
from the input formula permits an equisatisfiability claim only if any
previous interpretation of that name may be changed.
For example, on \(D=\{0,1\}\), \(P=\{1\}\), a declared \(c=0\)
cannot witness \(\exists xP(x)\), although \(c\) does not occur there.
Witness selection below uses ordinary set-theoretic choice for the
families of nonempty subsets of powers of \(D\).
\end{hbdefinition}

\begin{hbdefinition}{def:p12-opt-polarity}{Occurrence polarity}
The root has polarity \(+1\). Conjunction, disjunction, quantifiers,
and implication's consequent preserve polarity; negation and
implication's antecedent reverse it. Either child of a biconditional
has polarity \(0\), and descendants of a zero-polarity occurrence
retain polarity \(0\). Write \(C[A]\) for replacement of one designated
occurrence, preserving its surrounding binders.
\end{hbdefinition}

The exact contextual implication theorem, including every quantifier
case, is Theorem~\ref{thm:p12-small-context}. The closure and
rectification interfaces are Theorems~\ref{thm:p12-small-sat-closure}
and~\ref{thm:p12-small-rectification}. We use their all-assignment
conclusions when one newly interpreted symbol occurs under several
surrounding quantifiers.

\section{Formula renaming and exact clause costs}
\begin{hbdefinition}{def:p12-opt-renaming}{A polarity-sensitive name}
For an occurrence \(A=F|_\pi\), enumerate its distinct free variables
as \(\bar y\), including variables bound outside this occurrence.
Choose a fresh predicate \(R\) of arity \(|\bar y|\).
Replace the occurrence by \(R(\bar y)\) and conjoin the sentence
\[
 D_\pi=
 \begin{cases}
 \forall\bar y\,(R(\bar y)\to A),&\operatorname{pol}(F,\pi)=+1,\\
 \forall\bar y\,(A\to R(\bar y)),&\operatorname{pol}(F,\pi)=-1,\\
 \forall\bar y\,(R(\bar y)\leftrightarrow A),&
                  \operatorname{pol}(F,\pi)=0.
 \end{cases}
\]
Bound variables in the copied definition are renamed as necessary.
Zero free variables give a nullary predicate and an empty universal
prefix.
\end{hbdefinition}

\begin{hbtheorem}{thm:p12-opt-renaming}{proved=annex:p12.opt-renaming}{Naming preserves models by expansion and reduct}
Every model of \(F\) expands to a model of
\(F[\pi/R(\bar y)]\land D_\pi\); every model of the latter has a
\(\Sigma\)-reduct satisfying \(F\). Thus naming preserves satisfiability.
It does not assert equivalence for every arbitrary interpretation of \(R\).
\end{hbtheorem}

\begin{hbdefinition}{def:p12-opt-counts}{Counts before reduction}
For formulas without truth constants define \(p(A)\) to be the number
of clauses in the unreduced, polarity-sensitive standard translation;
write \(q(A)=p(\neg A)\). Count repeated clauses separately. No
tautology deletion, subsumption, equality simplification, or sharing
is included. With \(p_i=p(A_i)\) and \(q_i=q(A_i)\), the full recursion is
\[
\begin{array}{c|cc}
A&p(A)&q(A)\\ \hline
\text{atom}&1&1\\
\neg A_1&q_1&p_1\\
A_1\land A_2&p_1+p_2&q_1q_2\\
A_1\lor A_2&p_1p_2&q_1+q_2\\
A_1\to A_2&q_1p_2&p_1+q_2\\
A_1\leftrightarrow A_2&p_1q_2+q_1p_2&p_1p_2+q_1q_2\\
\forall xA_1,\ \exists xA_1&p_1&q_1
\end{array}
\]
First simplify constants; if the result is \(\top\) or \(\bot\),
return respectively the empty clause set or the set containing the
empty clause. This avoids applying the strictly positive count
identities to zero-clause cases.
\end{hbdefinition}

\begin{hbtheorem}{thm:p12-opt-counts}{proved=annex:p12.opt-counts}{Correct counts and finite arithmetic}
The recursion computes those clause counts. For a syntax tree of
\(N\) nodes, both counts are at most \(2^N\).
Thus their binary representations have at most \(N+1\) bits.
Dynamic programming computes them in polynomial bit complexity;
exponential numerical magnitude alone does not imply exponential
time to represent or calculate them.
\end{hbtheorem}

\begin{hbdefinition}{def:p12-opt-coefficients}{Coefficients of one occurrence}
For a fixed root formula \(F\), put \((a_\epsilon,b_\epsilon)=(1,0)\).
At a parent occurrence \(\tau\) with coefficients \(a,b\), the
coefficients at its child are as follows. The other child is \(A_j\).
\[
\begin{array}{c|c|cc}
\text{parent}&\text{child}&a_{\tau i}&b_{\tau i}\\ \hline
A_1\land A_2&i&a&bq_j\\
A_1\lor A_2&i&ap_j&b\\
A_1\to A_2&1&b&ap_2\\
A_1\to A_2&2&aq_1&b\\
A_1\leftrightarrow A_2&i&aq_j+bp_j&ap_j+bq_j\\
\neg A_1&1&b&a\\
\forall xA_1,\ \exists xA_1&1&a&b
\end{array}
\]
Counts are propagated from leaves to root; coefficients from root to
leaves. These directions refer to the syntax tree.
\end{hbdefinition}

\begin{hbtheorem}{thm:p12-opt-coefficients}{proved=annex:p12.opt-coefficients}{Exact one-occurrence saving test}
At any occurrence \(\pi\) there is a nonnegative integer \(c_\pi\)
independent of the formula inserted there such that
\[
 p(F[\pi/B])=a_\pi p(B)+b_\pi q(B)+c_\pi.
\]
For positive polarity \(a_\pi\ge1,b_\pi=0\); for negative polarity
\(a_\pi=0,b_\pi\ge1\); for zero polarity both are at least one.
Renaming \(A=F|_\pi\) does not increase the unreduced count exactly when
\[
\begin{array}{ll}
(a_\pi-1)(p(A)-1)\ge1,&\text{positive polarity},\\
(b_\pi-1)(q(A)-1)\ge1,&\text{negative polarity},\\
(a_\pi-1)(p(A)-1)+(b_\pi-1)(q(A)-1)\ge2,
 &\text{zero polarity}.
\end{array}
\]
The inequalities allow equality of the old and new counts.
\end{hbtheorem}

\begin{hbmethod}{def:p12-opt-capped}{Complete bounded-arithmetic implementation}
Represent a nonnegative integer \(v\) by
\(\langle v\rangle_3=\min(v,3)\). In both preceding tables replace
addition and multiplication by addition and multiplication followed
by this cap. The possible values are \(0,1,2,3\), with \(3\)
meaning at least three. Counts never equal zero; coefficients can.
For a positive-polarity test, both capped inputs must exceed one.
For a zero-polarity test, compute each product after subtracting one
from its two positive capped inputs, cap that product at two, and
test whether the sum is at least two.

For comparison, the complete threshold formulas for positive integers
\(u,v\) are
\[
\begin{array}{c|l}
u+v>2&u>1\ \text{or}\ v>1\\
uv>2&u>2\ \text{or}\ v>2\ \text{or}\ (u>1\ \text{and}\ v>1)\\
u+v>3&u>2\ \text{or}\ v>2\ \text{or}\ (u>1\ \text{and}\ v>1)\\
uv>3&u>3\ \text{or}\ v>3\ \text{or}\ (u>1\ \text{and}\ v>1).
\end{array}
\]
The last two lines explicitly include \(u=v=2\).
For a coefficient product with \(a\ge0,p\ge1\),
\(ap>1\) means \(a>1\) or \((a=1\ \text{and}\ p>1)\);
zero coefficients cannot be dropped from this test.
\end{hbmethod}

\begin{hbmethod}{def:p12-opt-schedule}{A finite top-down naming schedule}
Start with the rectified, constant-free input tree and give each original
nonliteral occurrence a distinct marker. Visit these markers in a fixed
preorder: a parent precedes its children and the first child precedes the
second. Keep the working formula as a tree with parent links and retain
each defining sentence as a separate conjunct of that formula.

At a marker, recompute the counts, coefficients and polarity in the
\emph{current whole formula}. If the corresponding nonincrease test
fails, remove that marker and proceed. If it succeeds, perform the
polarity-sensitive replacement of
Definition~\ref{def:p12-opt-renaming}, append its defining sentence,
and remove the visited marker. The copied subformula in the definition
inherits its descendants' still-unvisited markers; it replaces, rather
than duplicates, their former positions in the working tree. Rename
copied binders to avoid capture, retaining those occurrence markers.
The new name, definition connective and universal parameter prefix
receive no markers. Do not expand the defining biconditionals during
this traversal.

Stop after the original finite marker list has been exhausted, then
apply the standard clause translation to the resulting formula.
Each accepted step preserves satisfiability and does not increase the
stated unreduced clause count. This schedule fixes the termination and
cost interface; it makes no global optimality claim. Shared-instance
naming is a separate permitted operation with the conditions in
Theorem~\ref{thm:p12-opt-sharing}.
\end{hbmethod}

\begin{hbtheorem}{thm:p12-opt-capped}{proved=annex:p12.opt-capped}{Decision cost and controlled complete renaming}
The capped recursion decides every one-occurrence nonincrease test
exactly, in linear time and space in a fixed formula tree under
constant-cost operations on its finite value set.
The schedule in Definition~\ref{def:p12-opt-schedule} visits at most
\(N\) markers, preserves satisfiability, and gives a polynomial
algorithm. Its final unreduced clause count is no greater than the
original count.

A concrete complete naming algorithm names every nonliteral
occurrence once, using a full equivalence whose right-hand side
contains only its immediate connective and the names of its children.
Assert the root name, or the root literal if the input is already a
literal. Treat a quantifier node by its corresponding
quantified child name. Standard conversion of these shallow definitions
produces \(O(N)\) clauses. With all free variables carried as arguments
and ordinary inner Skolemization, its explicit symbol output is
\(O(N^3)\) in the original first-order tree size, including term sizes.
\end{hbtheorem}

\begin{hbexample}{def:p12-opt-renaming-cases}{Single and shared occurrences}
If \(p(A)=n,p(B)=m\), then \(A\lor\forall xB\) has \(nm\)
unreduced clauses, whereas naming \(B\), whose only free variable is
\(x\), gives \(n+m\). The inequality is \((n-1)(m-1)\ge1\).
For \(n=m=2\) the cost ties at four; for \(n=2,m=3\) it decreases
from six to five.

Each \(P_i\lor(Q_1\land Q_2)\) separately has two clauses and a
private name would give three. In the conjunction of \(n\) such
formulas, one shared positive definition
\(R\to(Q_1\land Q_2)\) instead gives the \(n+2\) clauses
\[
 P_i\lor R\ (1\le i\le n),\qquad
 \neg R\lor Q_1,\quad\neg R\lor Q_2.
\]
This beats \(2n\) exactly for \(n>2\).
For instances \(Q_1(a_i)\land Q_2(a_i)\), a sufficient common
template is \(Q_1(x)\land Q_2(y)\); use
\(\forall x,y(R(x,y)\to(Q_1(x)\land Q_2(y)))\) and replace the
instances by \(R(a_i,a_i)\).
\end{hbexample}

\begin{hbtheorem}{thm:p12-opt-sharing}{proved=annex:p12.opt-sharing}{Shared instances and syntactic matching}
Suppose finitely many nonoverlapping occurrences are capture-free
instances \(A\sigma_i\) of a template with free-variable list \(\bar y\).
One definition \(R(\bar y)\to A\) suffices if all occurrences are
positive; use \(A\to R(\bar y)\) if all are negative, and the full
equivalence if signs mix or any occurrence has zero polarity.
Replacing each by \(R(\bar y)\sigma_i\) preserves satisfiability
by expansion and reduct.
Ordinary syntactic matching of finite trees, with binders already
aligned and only free variables substitutable, is decidable in time
linear in the total pattern and target tree size.
\end{hbtheorem}

In the preceding propositional example, resolution of \(P_i\lor R\)
with each \(\neg R\lor Q_j\) reconstructs \(P_i\lor Q_j\).
Thus later inference can undo the syntactic saving. Restricting
resolution on auxiliary names requires an explicitly specified ordered
calculus; the present transformation does not establish that calculus.
Likewise, a left-associated chain containing only \(P,Q\) and
biconditionals has at most four assignments to its old atoms.
Repeated naming can introduce linearly many new atoms, hence
exponentially many possible assignments to those new atoms before
constraints are applied. After condensation has removed repeated
literals, a pairwise nonsubsuming, nontautological two-atom clause set
has at most four clauses: with no units only
the four binary clauses remain; one unit subsumes two of these;
two units allow at most one binary clause unless they are opposite
units, in which case they subsume all binary clauses; three or
four units subsume every binary clause. An empty clause subsumes
everything. Unreduced count alone misses those simplifications.
The condensation qualification matters when clauses are multisets.
For any \(r\ge1\), the \(r+1\) clauses containing \(i\) copies of
\(P\) and \(r-i\) copies of \(Q\), for \(0\le i\le r\), are
pairwise incomparable under multiset inclusion and contain no
tautology. Removing their repeated literals restores the ordinary
two-atom simplifications.

\section{Witness functions with fewer restrictions}
\begin{hbdefinition}{def:p12-opt-skolem-modes}{Inner, outer, and optimized steps}
Let \(F\) be a rectified sentence in negation normal form.
For an existential occurrence \(\exists xA\), inner Skolemization
uses \(f(\operatorname{FV}(\exists xA))\); outer Skolemization uses
the enclosing universally bound variables. The outer list may
include variables absent from this occurrence. A convenient
construction eliminates enclosing existential binders first, but
the next theorem justifies the source's arbitrary single-step
choice as well.

For optimized Skolemization choose an occurrence
\[
 E=\exists x_1\cdots x_k(A\land B),\qquad
 \bar y=\operatorname{FV}(E),
\]
and require the actual semantic premise
\(F\models\forall\bar y\exists\bar x A\).
With fresh functions \(f_1,\ldots,f_k\) of arity \(|\bar y|\),
write \(\sigma=\{x_i\mapsto f_i(\bar y):1\le i\le k\}\).
The transformed sentence is
\[
  (\forall\bar y\,A\sigma)\land F[\pi/B\sigma].
\]
The premise is not a syntactic test and must be proved or supplied
by the problem's retained theory. A search for it may require a
resource bound; failure to find it does not authorize the step.
\end{hbdefinition}

\begin{hbtheorem}{thm:p12-opt-standard-skolem}{proved=annex:p12.opt-standard-skolem}{The full inner and outer single-step domains}
Let \(\exists xA\) be any occurrence in a closed rectified NNF
sentence \(F\). An inner step using exactly its free-variable list,
or an outer step using exactly the enclosing universal-variable
list, preserves satisfiability by same-domain expansion and reduct.
This includes an occurrence beneath other existential binders.
For inner Skolemization the expansion can preserve equivalence
of the existential subformula and its replacement at all parameter
assignments. An outer step need not have this stronger local
equivalence; its guarantee concerns the whole sentence.
\end{hbtheorem}

\begin{hbtheorem}{thm:p12-opt-optimized-skolem}{proved=annex:p12.opt-optimized-skolem}{Optimized Skolemization}
Under the preceding conditions every model of \(F\) has an expansion
satisfying the transformed sentence, and every model of the latter
has a reduct satisfying \(F\).
\end{hbtheorem}

\begin{hbdefinition}{def:p12-opt-strong-skolem}{Free-variable splitting and the strong step}
For \(E=\exists\bar x(A_1\land\cdots\land A_n)\), \(n\ge1\),
with \(k\ge1\) distinct existential variables, fix an ordering
of the finitely many free variables. List the following sets in
that order to obtain pairwise disjoint variable lists \(Z_i\):
\[
 Z_1=\operatorname{FV}(A_1)\setminus\{\bar x\},\qquad
 Z_i=(\operatorname{FV}(A_i)\setminus\{\bar x\})
                 \setminus\bigcup_{j<i}Z_j.
\]
Their concatenation \(Z\) lists exactly \(\operatorname{FV}(E)\).
Empty lists are allowed. Remove a conjunct independent of all
\(\bar x\) from the existential scope if convenient; the theorem
also works without that preprocessing.
Let \(f_1,\ldots,f_k\) be fresh functions of arity \(|Z|\).
For each conjunct use fresh universally bound lists
\(W_{i+1},\ldots,W_n\) of the respective lengths of \(Z_{i+1},\ldots,Z_n\).
Define
\[
 \Xi=\bigwedge_{i=1}^n
 \forall W_{i+1}\cdots W_n\,
 A_i\{x_j\mapsto
 f_j(Z_1,\ldots,Z_i,W_{i+1},\ldots,W_n):1\le j\le k\}.
\]
Different conjuncts have disjoint new binders. The strong step replaces
the positive occurrence \(E\) in the NNF sentence by \(\Xi\).
\end{hbdefinition}

\begin{hbdefinition}{def:p12-opt-selectors}{Total selectors}
For a nonempty subset \(S\subseteq D^k\), a selector in this
chapter is a surjection \(h:D^k\to S\), viewed also as a map into
\(D^k\). It can be chosen as a retraction:
fix \(s_0\in S\), set \(h(t)=t\) on \(S\), and \(h(t)=s_0\)
elsewhere. The projection \(\operatorname{pr}_j\) returns the \(j\)-th
coordinate. There is no selector onto the empty set.
\end{hbdefinition}

\begin{hbtheorem}{thm:p12-opt-selector}{proved=annex:p12.opt-selector}{Uniform selector construction}
Fix a structure and a formula \(A(\bar x,\bar y)\) with no other
free variables. There are total functions \(h_j(\bar y,\bar w)\)
such that, at every parameter value \(\bar b\) where
\(\exists\bar xA\) holds, all values \(h(\bar b,\bar w)\) satisfy
\(A\), and every satisfying tuple occurs as such a value.
Consequently, for any additional formula \(B\), pointwise
\[
 \exists\bar x(A\land B)\quad\Longrightarrow\quad
 [\forall\bar w\,A(\bar h(\bar y,\bar w),\bar y)]
 \land\exists\bar v\,B\{\bar x\mapsto\bar h(\bar y,\bar v)\}.
\]
The functions are defined even at empty fibers; no truth of the
universal conjunct is claimed there.
\end{hbtheorem}

\begin{hbtheorem}{thm:p12-opt-strong-skolem}{proved=annex:p12.opt-strong-skolem}{Strong Skolemization at full first-order scope}
Every structure has an expansion for the functions in \(\Xi\)
such that \(E\) implies \(\Xi\) at every assignment to \(Z\).
In every expansion, \(\Xi\) implies \(E\).
It follows that replacing any existential-conjunction occurrence in a
rectified NNF sentence by \(\Xi\) preserves satisfiability by expansion
and reduct. Instantiating each \(W_j\) by \(Z_j\) obtains the
corresponding ordinary inner-Skolemized conjuncts.
The annex gives both the selector-composition route and a direct
prefix-witness construction.
\end{hbtheorem}

\begin{hbexample}{def:p12-opt-skolem-comparisons}{The two source comparisons}
In these examples \(R\) denotes a binary relation in the original
signature, rather than an auxiliary name. Consider
\[
 \forall x,y,z\,
 ((R(x,y)\land R(x,z))\to
             \exists u(R(y,u)\land R(z,u))).
\]
Inner Skolemization gives
\(\neg R(x,y)\lor\neg R(x,z)\lor R(y,f(y,z))\)
and the same first two literals followed by \(R(z,f(y,z))\).
If the whole theory proves \(\forall y\exists uR(y,u)\),
optimized Skolemization gives the unit \(R(y,f(y,z))\)
and the second original clause. The added unused parameter \(z\)
does not affect the premise in a nonempty domain.
Strong Skolemization instead replaces the first positive literal by
\(R(y,f(y,w))\). Its first clause is equivalent, under universal
closure, to
\(\neg R(x,y)\lor R(y,f(y,w))\):
instantiate \(z=y\) for one direction; the shorter clause implies
the longer in the other. This is the source's condensation.

Next take
\[
 \forall x,y(R(x,y)\to\exists z(P(z)\land Q(x,z)\land S(x,y,z)))
\]
and retain a proof that \(R\) is nonempty. The splitting is
\((\varnothing,(x),(y))\).
The standard clauses have \(\neg R(x,y)\) followed respectively by
\(P(f(x,y)),Q(x,f(x,y)),S(x,y,f(x,y))\).
Strong Skolemization changes the first two positive literals to
\(P(f(v,w))\) and \(Q(x,f(x,w))\).
Nonemptiness of \(R\) then implies the universal unit \(P(f(v,w))\).
Optimized Skolemization can also obtain the unit \(P(f(x,y))\),
since the original sentence and nonemptiness imply \(\exists zP(z)\);
its remaining \(Q\)-clause retains \(f(x,y)\).
Thus a globally proved existence condition and a local dependency
change offer different simplifications. Neither example proves a
universal ordering of the two methods' search costs.
\end{hbexample}

\begin{hbexample}{def:p12-opt-skolem-failures}{Necessary witness qualifications}
The optimized premise cannot be omitted. Let \(F=Q\lor
\exists x(P(x)\land P(x))\), on a one-element domain with \(Q\)
true and \(P\) empty. Then \(F\) is true, but pulling out the unit
\(P(f)\) is impossible in every expansion of this structure.
For dependence, \(\forall y\exists x(x=y)\) is true on a two-element
domain; replacing \(x\) by one constant makes it false there.
For context polarity, \(\neg\exists xP(x)\) must first be put in
NNF. Replacing its negative existential occurrence by \(P(c)\)
would admit a structure where \(P\) is nonempty but \(c\notin P\).
\end{hbexample}

\section{Equality and atom definitions}
\begin{hbtheorem}{thm:p12-opt-equality}{proved=annex:p12.opt-equality}{Identity-based simplification}
Interpret \(=\) as identity, let \(x\notin\operatorname{vars}(t)\),
and perform substitution without capture. In every structure and
assignment,
\[
 \exists x(x=t\land A)\ \Leftrightarrow\ A[x/t],
 \qquad
 \forall x(x=t\to A)\ \Leftrightarrow\ A[x/t].
\]
Also \(t=t\) is true. If \(y\notin\operatorname{vars}(t)\),
\(\forall y(t\ne y)\) is false and \(\exists y(t=y)\) is true.
On domains with more than one element, \(\forall y(t=y)\) is false
and \(\exists y(t\ne y)\) true; on singleton domains the respective
values are true and false.
Each replacement is valid in an arbitrary context on the stated
class of structures. For a background-theory application that class
must remain fixed on both sides.
\end{hbtheorem}

The quantifier-eliminating equality replacements decrease the number
of quantifiers even when substituting \(t\) enlarges terms. The
cardinality replacements decrease formula size. All are terminating
when used alone or with the standard size-decreasing Boolean rules.
For example
\(\forall x,y,z(x=y\to f(x,z)=f(y,z))\) becomes
\(\forall y,z(f(y,z)=f(y,z))\), then \(\top\).
Successively eliminate antecedent equalities to handle every
coordinate of function congruence; predicate congruence reduces to
\(P(\bar y)\to P(\bar y)\). Reflexivity, symmetry and transitivity
reduce in the same way, with symmetry becoming \(y=y\) after its
antecedent substitution and transitivity reducing its two equalities
in succession.
For a finite relevant signature there are finitely many such
congruence axioms. An arbitrary infinite signature instead has
a congruence scheme; the finite count must not be asserted for
all of its symbols at once.

A retained conjunct \(s\ne t\) certifies at least two elements for
other replacements. A statement merely about models of the old
whole formula is insufficient: the inconsistent formula
\((\forall y\,c=y)\land(\exists y\,c\ne y)\) vacuously implies
singleton cardinality, but replacing its first conjunct by \(\top\)
creates models on larger domains.
Nor can \(x\notin\operatorname{vars}(t)\) be removed: on two elements
with \(f\) interchanging them, \(\exists x(x=f(x)\land\top)\)
is false, whereas substitution into \(\top\) gives true.

\begin{hbdefinition}{def:p12-opt-atom-definition}{A retained atom definition}
Let \(\Theta\) be a retained sentence implying
\[
 \forall\bar x(P(\bar t)\leftrightarrow A),
 \qquad \{\bar x\}=\operatorname{FV}(P(\bar t)),\quad
 \operatorname{FV}(A)\subseteq\{\bar x\}.
\]
An occurrence \(P(\bar s)=P(\bar t)\sigma\), with capture-free
matching substitution \(\sigma\), may be replaced by \(A\sigma\).
The retained context is \(\Theta\land C[P(\bar s)]\).
For a conditional definition
\(\forall\bar x(H\to(P(\bar t)\leftrightarrow A))\),
also require \(\operatorname{FV}(H)\subseteq\{\bar x\}\), so that
the displayed definition is a sentence.
The guard \(H\sigma\) must hold at every assignment relevant to the
replacement. One sufficient explicit interface is
\(\Theta\models\forall\operatorname{FV}(H\sigma)\,H\sigma\).
Another is an occurrence inside \(H\sigma\to C[P(\bar s)]\),
provided no binder between the guard and the occurrence binds
a free variable of \(H\sigma\). Outside that guard the whole
implication is already true.
\end{hbdefinition}

\begin{hbtheorem}{thm:p12-opt-atom-expansion}{proved=annex:p12.opt-atom-expansion}{Expansion, elimination and termination}
These replacements preserve truth under the retained conditions at
every polarity. For one unconditional definition whose head is
\(P(x_1,\ldots,x_n)\) with distinct variables and whose body contains
no \(P\), expansion eliminates all targeted \(P\)-occurrences in
finitely many steps.
If a theory consists of this definition and formulas using \(P\),
fully expand \(P\) in the other formulas and delete the defining
sentence. The resulting theory over \(\Sigma\setminus\{P\}\)
has exactly the reducts of the original models; every such reduct
has a unique expansion for \(P\).
\end{hbtheorem}

Expansion can increase size. A recursive definition
\(\forall x(P(x)\leftrightarrow P(f(x)))\) allows the endless
sequence \(P(a),P(f(a)),P(f(f(a))),\ldots\).
Several individually nonrecursive definitions can also cycle,
as \(P\leftrightarrow Q\), \(Q\leftrightarrow P\).
With a guarded definition, \(H\to(P\leftrightarrow Q)\),
a valuation with \(H\) false, \(P\) true and \(Q\) false satisfies
the definition but refutes an unguarded replacement of \(P\) by \(Q\).
Naming to reduce a predicted clause count and expanding a domain
definition therefore serve different purposes.

\section{Decision diagrams}
\begin{hbdefinition}{def:p12-opt-bdd}{Ordered reduced Boolean decision diagrams}
For a finite list of propositional variables
\(v_1<\cdots<v_m\), a decision diagram is a finite rooted acyclic
graph whose vertices are all reachable from its root, with terminal
labels in \(\{0,1\}\) and internal nodes \((v_i,L,H)\).
Either terminal may be absent when it is not reachable.
Follow \(L\) when \(v_i=0\), \(H\) when \(v_i=1\).
Along every path variable indices strictly increase.
Reduction identifies equal terminals, merges nodes with identical
triples, and replaces a node with \(L=H\) by its child.
Discard any vertices made unreachable by a replacement.
The result is reduced when neither operation changes it.
\end{hbdefinition}

\begin{hbtheorem}{thm:p12-opt-bdd}{proved=annex:p12.opt-bdd}{Shannon splitting, reduction and uniqueness}
For any designated subformula \(A\) of a propositional formula \(F\),
\[
 F\ \Leftrightarrow\
 (A\to F[\pi/\top])\land(\neg A\to F[\pi/\bot]).
\]
Full splitting by a fixed variable order followed by reduction
terminates, preserves truth, and yields a unique reduced graph
up to graph isomorphism respecting labels and edges.
Truth and falsity yield the respective single terminals.
For each root-to-\(0\) path, form the clause that forbids its partial
assignment: negate each high-branch variable and leave each
low-branch variable positive. The conjunction of these clauses is
equivalent to the represented Boolean function.
\end{hbtheorem}

For the source example \((Q\leftrightarrow R)\leftrightarrow
(Q\leftrightarrow S)\), setting \(Q\) true gives
\(R\leftrightarrow S\); setting it false gives
\((\neg R)\leftrightarrow(\neg S)\), the same function.
Delete the \(Q\)-node. The two paths to \(0\) are
\((R,S)=(1,0)\) and \((0,1)\); their clauses are
\(\neg R\lor S\) and \(R\lor\neg S\).

At a fixed first-order scope one may regard maximal quantified
subformulas and atoms as Boolean tests. Every Boolean identity
obtained this way remains sound after those tests are replaced
by their formulas. It is not a complete test of first-order
equivalence: \(\forall xP(x)\to P(c)\) is valid over nonempty
first-order structures, but its Boolean skeleton \(U\to V\)
has a falsifying valuation. Moving a test across a quantifier
requires its own free-variable conditions.
Full splitting has as many as \(2^{m+1}-1\) nodes before sharing;
this construction gives no general polynomial bound on diagram
size or on the number of root-to-false paths.

\section{Historical and implementation scope}
\label{sec:p12-opt-history}
The historical claims here belong to the 2001 chapter's
Sections 7--8 \cite[pp.~363--365]{NonnengartWeidenbach2001SmallClauses}.
It credits polarity-sensitive equivalence elimination to the 1980
work on Andrews's Challenge Problem 1; the introduction of new
propositional names to Tseitin; and formula-renaming developments
to Eder, Plaisted--Greenbaum, and Boy de la Tour.
The selected criterion follows Boy de la Tour's count objective,
presented using positions after Rock and a bounded-arithmetic
implementation after Nonnengart--Rock--Weidenbach.

The chapter reports a separate Boy de la Tour optimality result:
for formulas without biconditionals, a top-down traversal using the
count criterion obtains an optimum among the considered renamings;
strict saving rather than nonincrease still suffices in that setting.
This is an attribution of that narrower historical result, whose
original optimality proof is not supplied by the selected chapter.
It is not a theorem here about arbitrary logical equivalence,
shared-instance matching, redundancy removal, or proof-search time.
The chapter also reports a quadratic symbol bound for the
Plaisted--Greenbaum approach, and a cubic bound for its own
approach. The explicit cubic construction proved above declares
its data and output model; these historical algorithm variants
must not be identified merely because their bounds have the same order.

The selected source directs readers to Skolem and later textbook
treatments for standard Skolemization; Egly for antiprenexing;
Ohlbach--Weidenbach, Nonnengart, and the 1998 joint paper for the
advanced witness methods; Baaz--Leitsch for proof-complexity
questions; and Goubault and Goubault--Posegga for first-order
BDD simplification and combinations with Skolemization.
Reuse of witness functions and introduced predicates is explicitly
outside its treatment. Citing these works here records the
source's referral, not acquisition of each original presentation.

Its implementation discussion concerns the then-described SPASS
work, with a reference to SPASS version 1.0.0. It recommends a
top-down traversal with efficient access to children and parent,
and storing formula size as a prefilter for syntactic instance
matching. Structural matching is the relation
\(A\sigma=B\) for some substitution \(\sigma\),
under the chapter's restricted substitution convention; it is
distinct from matching modulo associativity, commutativity,
or logical implication. For capture-safe manuscript use, the
binder-aligned interface of
Theorem~\ref{thm:p12-opt-sharing} specifies the admissible substitutions.

The source reports favorable experience and refers to the 1998 paper
for experiments, but these selected pages give no benchmark counts,
timing table, hardware specification, confidence intervals or
replication data. They support a historical report of the authors'
experience, not a numerical speedup or present-day ranking.
Their suggestion to combine optimized and strong Skolemization
remains a heuristic proposal illustrated by the two worked cases.

\begin{proofboundary}{def:p12-opt-remaining-originals}{Exact remaining source obligations}
The selected chapter's naming, optimized and strong witness
constructions, and three simplification methods have explicit
proofs above and in the annex. Separate original-source obligations
remain for Boy de la Tour's global optimality argument, the distinct
Plaisted--Greenbaum size analysis, the cited proof-length speedups,
the 1998 experimental comparison, and the first-order BDD algorithms
and witness-reuse methods only referred to here. The standard
rectification, NNF, closure and clause-generation pipeline is the
companion chapter on standard small-clause translation.
\end{proofboundary}
